\chapter{Conclusioni e sviluppi futuri}
\label{cap:conclusioni}

\notaredazionale{Obiettivo: 3 pagine. Riassumere il contributo del lavoro rispetto agli obiettivi enunciati nel
Capitolo \ref{cap:introduzione}, senza ripetere in dettaglio quanto già argomentato nei capitoli precedenti.

Limiti noti da richiamare (vedi \texttt{REPORT\_ARCHITETTURA.md} \S7 per l'elenco completo e aggiornato):
\begin{itemize}
    \item nessuna verifica automatica di fedeltà numerica prima della consegna della risposta;
    \item il rifiuto su domande fuori tema/senza informazione resta debole, indipendentemente dal modello;
    \item nessuna gestione della validità temporale dei documenti;
    \item la fusione gerarchica copre solo la sintesi interna a una sezione/documento, non tra documenti diversi;
    \item recall debole della ricerca ibrida su domande lunghe e lessicalmente sovraccariche del nome del corso.
\end{itemize}

Sviluppi futuri concreti già individuati:
\begin{itemize}
    \item batching dinamico delle richieste di embedding/reranking, per servire più studenti davvero in concorrenza
      sulla GPU condivisa (alternativa più scalabile del lock attuale);
    \item autenticazione reale / integrazione con il sistema informativo studenti, per eliminare alla radice il
      problema della disambiguazione del corso (vedi anche Capitolo \ref{cap:stato-arte}, \S2.4);
    \item ri-esecuzione della valutazione sperimentale del Capitolo \ref{cap:validazione} dopo le correzioni più
      recenti, per avere numeri aggiornati;
    \item eventuale estensione della fusione gerarchica all'aggregazione di informazione sparsa su più documenti.
\end{itemize}}
